\ifdefined\gkver\else\def\gkver{student}\fi
\documentclass[\gkver]{gaokaozhenti}
\begin{document}

\chapter{2014年广东理综}
%% paperType: 理综物理部分

\begin{enumerate}
\item
%% number: 13
%% paperName: 2014年广东理综第 13 题 4 分
%% typeId: 1
%% type: 单选题
%% chapter: 直线运动
%% point: 运动图象
%% method: 无
%% score: 4
%% degree: 780
%% duplicateId: 0
%% body:
如图所示是物体做直线运动的 $v-t$ 图象，由图可知，该物体\xzanswer{B}
\onepicture[align=right, w=3.6cm]{figs/13.png}
\fourchoices[answer=B]
{第 $1\Us$ 内和第 $3\Us$ 内的运动方向相反}
{第 $3\Us$ 内和第 $4\Us$ 内的加速度相同}
{第 $1\Us$ 内和第 $4\Us$ 内的位移大小不等}
{$0\sim2\Us$ 内和 $0\sim4\Us$ 内的平均速度大小相等}
%% answer: B
%% memo:
\memoanswer{
第 $1\Us$ 内和第 $3\Us$ 内的速度均为正值，方向相同，选项 A 错误；$v-t$ 图象的斜率代表加速度，第 $3\Us$ 内和第 $4\Us$ 内斜率相同，所以加速度相同，选项 B 正确；图象与时间轴所围面积在数值上等于位移的大小，第 $1\Us$ 内的位移 $x_{1}=\frac{1}{2}\times1\times1\Um=0.5\Um$，第 $4\Us$ 内的位移 $x_{4}=-\frac{1}{2}\times1\times1\Um=-0.5\Um$，两段时间内位移大小相等，选项 C 错误；$0\sim2\Us$ 内的平均速度 $\overline{v}=\frac{x}{t}=\frac{1.5}{2}\Ums=0.75\Ums$，$0\sim4\Us$ 内的平均速度 $\overline{v}'=\frac{x'}{t'}=\frac{1.5}{4}\Ums=0.375\Ums$，选项 D 错误。
}

\item
%% number: 14
%% paperName: 2014年广东理综第 14 题 4 分
%% typeId: 1
%% type: 单选题
%% chapter: 力物体的平衡
%% point: 受力分析
%% method: 无
%% score: 4
%% degree: 650
%% duplicateId: 0
%% body:
如图所示，水平地面上堆放着原木，关于原木 P 在支撑点 M、N 处受力的方向，下列说法正确的是\xzanswer{A}
\onepicture[align=right, w=5.4cm]{figs/14.png}
\fourchoices[answer=A]
{M 处受到的支持力竖直向上}
{N 处受到的支持力竖直向上}
{M 处受到的摩擦力沿 MN 方向}
{N 处受到的摩擦力沿水平方向}
%% answer: A
%% memo:
\memoanswer{
原木 P 的支撑点 M、N 处弹力垂直于接触面，M 处的支持力竖直向上，N 处受到的支持力垂直 MN 斜向上，故 A 项正确、B 项错误；M 处受到的静摩擦力沿水平方向，C 项错误；N 处受到的静摩擦力沿 MN 方向，D 项错误。
}

\item
%% number: 15
%% paperName: 2014年广东理综第 15 题 4 分
%% typeId: 1
%% type: 单选题
%% chapter: 电磁感应
%% point: 楞次定律推广
%% method: 无
%% score: 4
%% degree: 650
%% duplicateId: 0
%% body:
如图所示，上下开口、内壁光滑的铜管 P 和塑料管 Q 竖直放置，小磁块先后在两管中从相同高度处由静止释放，并落至底部，则小磁块\xzanswer{C}
\onepicture[align=right, w=3.2cm]{figs/15.png}
\fourchoices[answer=C]
{在 P 和 Q 中都做自由落体运动}
{在两个下落过程中的机械能都守恒}
{在 P 中的下落时间比在 Q 中的长}
{落至底部时在 P 中的速度比在 Q 中的长}
%% answer: C
%% memo:
\memoanswer{
本题原卷及材料中未提供详解，待补充。
}

\item
%% number: 16
%% paperName: 2014年广东理综第 16 题 4 分
%% typeId: 1
%% type: 单选题
%% chapter: 功和能
%% point: 功能原理
%% method: 无
%% score: 4
%% degree: 650
%% duplicateId: 0
%% body:
如图是安装在列车车厢之间的摩擦缓冲器结构图，图中①和②为楔块，③和④为垫块，楔块与弹簧盒、垫块间均有摩擦，在车厢相互撞击时弹簧压缩过程中\xzanswer{B}
\onepicture[align=right, w=5.4cm]{figs/16.png}
\fourchoices[answer=B]
{缓冲器的机械能守恒}
{摩擦力做功消耗机械能}
{垫块的动能全部转化成内能}
{弹簧的弹性势能全部转化为动能}
%% answer: B
%% memo:
\memoanswer{
本题原卷及材料中未提供详解，待补充。
}

\item
%% number: 17
%% paperName: 2014年广东理综第 17 题 6 分
%% typeId: 2
%% type: 多选题
%% chapter: 气体性质
%% point: 热力学第一定律
%% method: 无
%% score: 6
%% degree: 650
%% duplicateId: 0
%% body:
用密封性好、充满气体的塑料袋包裹易碎品，如图所示，充气袋四周被挤压时，假设袋内气体与外界无热交换，则袋内气体\xzanswer{AC}
\onepicture[align=right, w=4.6cm]{figs/17.png}
\fourchoices[answer=AC]
{体积减小，内能增大}
{体积减小，压强减小}
{对外界做负功，内能增大}
{对外界做正功，压强减小}
%% answer: AC
%% memo:
\memoanswer{
本题原卷及材料中未提供详解，待补充。
}

\item
%% number: 18
%% paperName: 2014年广东理综第 18 题 6 分
%% typeId: 2
%% type: 多选题
%% chapter: 光学
%% point: 光电效应
%% method: 无
%% score: 6
%% degree: 650
%% duplicateId: 0
%% body:
在光电效应实验中，用频率为 $\nu$ 的光照射光电管阴极，发生了光电效应，下列说法正确的是\xzanswer{AD}
\fourchoices[answer=AD]
{增大入射光的强度，光电流增大}
{减小入射光的强度，光电效应现象消失}
{改变频率小于 $\nu$ 的光照射，一定不发生光电效应}
{改变频率大于 $\nu$ 的光照射，光电子的最大初动能变大}
%% answer: AD
%% memo:
\memoanswer{
本题原卷及材料中未提供详解，待补充。
}

\item
%% number: 19
%% paperName: 2014年广东理综第 19 题 6 分
%% typeId: 2
%% type: 多选题
%% chapter: 交流电
%% point: 变压器
%% method: 无
%% score: 6
%% degree: 760
%% duplicateId: 0
%% body:
如图所示的电路中，P 为滑动变阻器的滑片，保持理想变压器的输入电压 $U_{1}$ 不变，闭合电键 S，下列说法正确的是\xzanswer{BD}
\onepicture[align=right, w=4.8cm]{figs/19.png}
\fourchoices[answer=BD]
{P 向下滑动时，灯 L 变亮}
{P 向下滑动时，变压器的输出电压不变}
{P 向上滑动时，变压器的输入电流变小}
{P 向上滑动时，变压器的输出功率变大}
%% answer: BD
%% memo:
\memoanswer{
本题原卷及材料中未提供详解，待补充。
}

\item
%% number: 20
%% paperName: 2014年广东理综第 20 题 6 分
%% typeId: 2
%% type: 多选题
%% chapter: 电场
%% point: 带电粒子的平衡
%% method: 无
%% score: 6
%% degree: 650
%% duplicateId: 0
%% body:
如图所示，光滑绝缘的水平桌面上，固定着一个带电量为 $+Q$ 的小球 P，带电量分别为 $-q$ 和 $+2q$ 的小球 M 和 N，由绝缘细杆相连，静止在桌面上，P 与 M 相距 $L$，P、M 和 N 视为点电荷，下列说法正确的是\xzanswer{BD}
\onepicture[align=right, w=4.2cm]{figs/20.png}
\fourchoices[answer=BD]
{M 与 N 的距离大于 $L$}
{P、M 和 N 在同一直线上}
{在 P 产生的电场中，M、N 处的电势相同}
{M、N 及细杆组成的系统所受合外力为零}
%% answer: BD
%% memo:
\memoanswer{
本题原卷及材料中未提供详解，待补充。
}

\item
%% number: 21
%% paperName: 2014年广东理综第 21 题 4 分
%% typeId: 2
%% type: 多选题
%% chapter: 曲线运动
%% point: 人造地球卫星
%% method: 无
%% score: 4
%% degree: 450
%% duplicateId: 0
%% body:
如图所示，飞行器 P 绕某星球做匀速圆周运动，星球相对飞行器的角度为 $\theta$，下列说法正确的是\xzanswer{AC}
\onepicture[align=right, w=3.2cm]{figs/21.png}
\fourchoices[answer=AC]
{轨道半径越大，周期越长}
{轨道半径越大，速度越长}
{若测得周期和张角，可得到星球的平均密度}
{若测得周期和轨道半径，可得到星球的平均密度}
%% answer: AC
%% memo:
\memoanswer{
本题原卷及材料中未提供详解，待补充。
}

\item
%% number: 34
%% paperName: 2014年广东理综第 34 题 10 分
%% typeId: 4
%% type: 实验题
%% chapter: 功和能
%% point: 功能原理
%% method: 无
%% score: 10
%% degree: 650
%% duplicateId: 0
%% body:
\begin{enumerate}
	\item
	（8 分）某同学设计的可调电源电路如图（a）所示，$R_{0}$ 为保护电阻，P 为滑动变阻器的滑片，闭合电键 S。

	\twopicture[numA=（a）, numB=（b）, labelprefix={}, wA=3.5cm, wB=4.6cm]{figs/34a.png}{figs/34b.png}

	\begin{steps}
		\item
		用电压表测量 A、B 两端的电压；将电压表调零，选择 $0\sim3\UV$ 档，示数如图（b），电压值为\tkanswer[1.4cm]{1.30}$\UV$。
		\item
		在接通外电路之前，为了保证外电路的安全，滑片 P 应先置于\tkanswer[1.2cm]{A}端。
		\item
		要使输出电压 $U$ 变大，滑片 P 应向\tkanswer[1.2cm]{B}端移动。
		\item
		若电源电路中不接入 $R_{0}$，则在使用过程中，存在\tkanswer[1.6cm]{短路}的风险（填“断路”或“短路”）。
	\end{steps}
	\item
	（10 分）某同学根据机械能守恒定律，设计实验探究弹簧的弹性势能与压缩量的关系。

	\threepicture[numA=（a）, numB=（b）, numC=（c）, labelprefix={}, wA=2.2cm, wB=4.9cm, wC=4.2cm]{figs/34c.png}{figs/34d.png}{figs/34e.png}

	\begin{steps}
		\item
		如图（a），将轻质弹簧下端固定于铁架台，在上端的托盘中依次增加砝码，测得相应的弹簧长度，部分数据如下表，由数据算得劲度系数 $k=$\tkanswer[1.4cm]{50}$\UNm$（$g$ 取 $9.8\Umsq$）。

		\begin{center}
		\begin{tabular}{|c|c|c|c|}
		\hline
		砝码质量（$\Ug$） & 50 & 100 & 150 \\
		\hline
		弹簧长度（$\Ucm$） & 8.62 & 7.63 & 6.66 \\
		\hline
		\end{tabular}
		\end{center}

		\item
		取下弹簧，将其一端固定于气垫导轨左侧，如图（b）所示；调整导轨，使滑块自由滑动时，通过两个光电门的速度大小\tkanswer[1.6cm]{相等}。
		\item
		用滑块压缩弹簧，记录弹簧的压缩量 $x$；释放滑块，记录滑块脱离弹簧后的速度 $v$，释放滑块过程中，弹簧的弹性势能转化为\tkanswer[2.0cm]{滑块的动能}。
		\item
		重复③中的操作，得到 $v$ 与 $x$ 的关系如图（c）。由图可知，$v$ 与 $x$ 成\tkanswer[1.4cm]{正比}关系，由上述实验可得结论：对同一根弹簧，弹性势能与弹簧的\tkanswer[1.8cm]{压缩量}成正比。
	\end{steps}
\end{enumerate}
%% answer:
\jdanswer{
\begin{enumerate}
	\item
	①1.30，②A，③B，④短路
	\item
	①50，②相等，③滑块的动能，④正比，压缩量
\end{enumerate}
}
%% memo:
\memoanswer{
本题原卷及材料中未提供详解，待补充。
}

\item
%% number: 35
%% paperName: 2014年广东理综第 35 题 18 分
%% typeId: 6
%% type: 计算题
%% chapter: 运动和力
%% point: 动量守恒定律
%% method: 无
%% score: 18
%% degree: 450
%% duplicateId: 0
%% body:
如图所示的水平轨道中，AC 段的中点 B 的正上方有一探测器，C 处有一竖直挡板，物体 $P_{1}$ 沿轨道向右以速度 $v_{1}$ 与静止在 A 点的物体 $P_{2}$ 碰撞，并接合成复合体 P，以此碰撞时刻为计时零点，探测器只在 $t_{1}=2\Us$ 至 $t_{2}=4\Us$ 内工作，已知 $P_{1}$、$P_{2}$ 的质量都为 $m=1\Ukg$，P 与 AC 间的动摩擦因数为 $\mu=0.1$，AB 段长 $L=4\Um$，$g$ 取 $10\Umsq$，$P_{1}$、$P_{2}$ 和 P 均视为质点，P 与挡板的碰撞为弹性碰撞。
\begin{enumerate}
	\item
	若 $v_{1}=6\Ums$，求 $P_{1}$、$P_{2}$ 碰后瞬间的速度大小 $v$ 和碰撞损失的动能 $\Delta E$；
	\item
	若 P 与挡板碰后，能在探测器的工作时间内通过 B 点，求 $v_{1}$ 的取值范围和 P 向左经过 A 点时的最大动能 $E$。
\end{enumerate}
\onepicture[align=right, w=8cm]{figs/35.png}
%% answer:
\jdanswer{
\begin{enumerate}
	\item
	$v=3\Ums$，$\Delta E=9\UJ$
	\item
	$10\Ums\leq v_{1}\leq14\Ums$，$E=49\UJ$
\end{enumerate}
}
%% memo:
\memoanswer{
解：（1）$P_{1}$、$P_{2}$ 碰撞过程，动量守恒

$mv_{1}=2mv$ ①

解得 $v=\frac{v_{1}}{2}=3\Ums$ ②

碰撞损失的动能 $\Delta E=\frac{1}{2}mv_{1}^{2}-\frac{1}{2}(2m)v^{2}$ ③

解得 $\Delta E=9\UJ$ ④

（2）由于 P 与挡板的碰撞为弹性碰撞，故 P 在 AC 间等效为匀减速运动，设 P 在 AC 段加速度大小为 $a$，由运动学规律，得

$\mu(2m)g=2ma$ ⑤

$3L=vt-\frac{1}{2}at^{2}$ ⑥

$v_{2}=v-at$ ⑦

由①⑤⑥⑦解得 $v_{1}=\frac{t^{2}+24}{t}$，$v_{2}=\frac{24-t^{2}}{2t}$ ⑧

由于 $2\Us\leq t\leq4\Us$，所以解得 $v_{1}$ 的取值范围 $10\Ums\leq v_{1}\leq14\Ums$ ⑨

$v_{2}$ 的取值范围 $1\Ums\leq v_{2}\leq7\Ums$

所以当 $v_{2}=7\Ums$ 时，P 向左经过 A 点时有最大动能

$E=\frac{1}{2}(2m)v_{2}^{2}=49\UJ$ ⑩
}

\item
%% number: 36
%% paperName: 2014年广东理综第 36 题 18 分
%% typeId: 6
%% type: 计算题
%% chapter: 磁场
%% point: 带电粒子在磁场中的运动
%% method: 无
%% score: 18
%% degree: 250
%% duplicateId: 0
%% body:
如图所示，足够大的平行挡板 $A_{1}$、$A_{2}$ 竖直放置，间距 $6L$。两板间存在两个方向相反的匀强磁场区域Ⅰ和Ⅱ，以水平面 MN 为理想分界面，Ⅰ区的磁感应强度为 $B_{0}$，方向垂直纸面向外。$A_{1}$、$A_{2}$ 上各有位置正对的小孔 $S_{1}$、$S_{2}$，两孔与分界面 $MN$ 的距离均为 $L$。质量为 $m$、电量为 $+q$ 的粒子经宽度为 $d$ 的匀强电场由静止加速后，沿水平方向从 $S_{1}$ 进入Ⅰ区，并直接偏转到 MN 上的 P 点，再进入Ⅱ区，P 点与 $A_{1}$ 板的距离是 $L$ 的 $k$ 倍，不计重力，碰到挡板的粒子不予考虑。
\begin{enumerate}
	\item
	若 $k=1$，求匀强电场的电场强度 $E$；
	\item
	若 $2<k<3$，且粒子沿水平方向从 $S_{2}$ 射出，求出粒子在磁场中的速度大小 $v$ 与 $k$ 的关系式和Ⅱ区的磁感应强度 $B$ 与 $k$ 的关系式。
\end{enumerate}
\onepicture[align=right, w=6cm]{figs/36.png}
%% answer:
\jdanswer{
\begin{enumerate}
	\item
	$E=\frac{qB_{0}^{2}L^{2}}{2md}$
	\item
	$v=\frac{(k^{2}+1)qB_{0}L}{2m}$，$B=\frac{k}{3-k}B_{0}$
\end{enumerate}
}
%% memo:
\memoanswer{
解：（1）粒子在电场中，由动能定理有

$qEd=\frac{1}{2}mv^{2}-0$ ①

粒子在Ⅰ区洛伦兹力提供向心力

$qvB_{0}=m\frac{v^{2}}{r}$ ②

当 $k=1$ 时，由几何关系得

$r=L$ ③

由①②③解得

$E=\frac{qB_{0}^{2}L^{2}}{2md}$ ④

（2）由于 $2<k<3$ 时，由题意可知粒子在Ⅱ区只能发生一次偏转，由几何关系可知

$(r-L)^{2}+(kL)^{2}=r^{2}$ ⑤

解得 $r=\frac{k^{2}+1}{2}L$ ⑥

由②⑥解得 $v=\frac{(k^{2}+1)qB_{0}L}{2m}$ ⑦

粒子在Ⅱ区洛伦兹力提供向心力

$qvB=m\frac{v^{2}}{r_{1}}$ ⑧

由对称性及几何关系可知 $\frac{k}{3-k}=\frac{r}{r_{1}}$ ⑨

解得 $r_{1}=\frac{(3-k)(k^{2}+1)}{2k}L$ ⑩

由⑧⑩解得 $B=\frac{k}{3-k}B_{0}$
}

\end{enumerate}

\end{document}
